\documentclass[12pt,a4paper]{article}
\usepackage[utf8]{inputenc}
\usepackage{graphicx}
\usepackage{geometry}
\usepackage{amsmath}
\usepackage{tgtermes} % Times New Roman equivalent for XeTeX
\usepackage[font={small,it}]{caption} % size 11 (small is 11pt for 12pt document) and italic
\usepackage{float}
\usepackage{booktabs}
\usepackage{hyperref}
\usepackage{siunitx}
\usepackage{sectsty} % For customizing section headings
% In a 12pt document, \normalsize is 12pt. \bfseries makes it bold.
\sectionfont{\fontsize{12}{14}\bfseries\selectfont} 

\geometry{
  a4paper,
  left=25mm,
  right=25mm,
  top=25mm,
  bottom=25mm,
}

\begin{document}

% --- Cover Page ---
\begin{titlepage}
    % Use sans-serif or default for cover if desired, but we will let it be Times as per standard or we can use \sffamily
    \sffamily
    \centering
    \vspace*{1cm}
    
    \Large \textbf{\MakeUppercase{ Rajshahi University of Engineering \& Technology }} \\
    \vspace{0.5cm}
    \large \textbf{\MakeUppercase{ Department of Electrical and Electronic Engineering }} \\
    
    \vspace{2.5cm}
    \Large \textbf{LABORATORY REPORT} \\
    \vspace{1cm}
    
    \begin{tabular}{ll}
        \textbf{Course No:} & EEE 3102 \\
        \textbf{Course Title:} & Digital Signal Processing Sessional \\
    \end{tabular}
    
    \vspace{2cm}
    
    \textbf{Experiment No:} 03 \\
    \vspace{0.5cm}
    \textbf{Name of the Experiment:} \\
    \Large \textbf{ Test Experiment about triac and diac together } \\
    
    \vspace{2.5cm}
    
    \begin{minipage}[t]{0.4\textwidth}
        \begin{flushleft} \large
            \emph{Submitted by:}\\
            John Doe \\
            Roll: 1901000 \\
            Section: A, Group: 1
        \end{flushleft}
    \end{minipage}
    \hfill
    \begin{minipage}[t]{0.5\textwidth}
        \begin{flushright} \large
            \emph{Submitted to:} \\
            
            Dr. Jane Smith \\
            Professor \\
            \vspace{0.3cm}
            
            Mr. Bob Brown \\
            Lecturer \\
            
            
        \end{flushright}
    \end{minipage}

    \vfill
    \textbf{Date of Performance:} October 09, 2026 \\
    \textbf{Date of Submission:} October 16, 2026
    
\end{titlepage}

% --- Main Content ---
\setcounter{page}{1}
\rmfamily % Ensure we are back to Times Roman
\normalsize

\begin{center}
    \textbf{Experiment No:} 03 \\
    \vspace{0.2cm}
    \textbf{Name of the Experiment:} Test Experiment about triac and diac together
\end{center}
\vspace{0.5cm}


\section{Objectives}
\begin{itemize}

    \item To investigate the operational characteristics of a TRIAC and DIAC combination in an AC power control circuit.

    \item To determine the effect of the DIAC breakover voltage on the firing angle of the TRIAC gate trigger.

    \item To analyze the bidirectional conduction behavior of the TRIAC when triggered by a DIAC in both positive and negative half-cycles.

\end{itemize}


\section{Introduction}
Thyristor-based AC power control is a fundamental application in modern electrical engineering, utilizing semiconductor devices to regulate power to resistive and inductive loads. Central to these systems are the TRIAC (Triode for Alternating Current) and the DIAC (Diode AC switch). A TRIAC is a bidirectional switching device that can conduct current in both directions when triggered by a gate signal, effectively acting as a bidirectional SCR. In its off-state, it blocks voltage until a specific gate-to-cathode voltage is applied. However, manually generating precise gate pulses in an AC cycle is impractical; therefore, a DIAC is commonly employed as a bidirectional trigger source. The DIAC is a two-terminal, symmetric, latching device that remains non-conductive until the voltage across it exceeds its breakover voltage, typically ranging from 30 V to 80 V depending on the specific part. Once this threshold is reached, the DIAC latches on and conducts a sharp pulse of current, which is delivered to the TRIAC gate, forcing the TRIAC into conduction.

The interaction between the DIAC and TRIAC allows for phase control of the load voltage. By varying the RC time constant in the trigger circuit, the point in the AC cycle at which the DIAC breaks over can be adjusted. This delay, known as the firing angle (α), determines how much of the sine wave is delivered to the load. When α is small, the TRIAC conducts for a larger portion of the cycle, resulting in higher power transfer; when α is large, conduction is shortened, reducing power. The symmetrical nature of the DIAC ensures that the TRIAC is triggered at the same phase angle in both the positive and negative half-cycles, providing balanced AC control without rectification. This principle is widely used in lamp dimmers, fan speed controllers, and motor soft-starters.



\section{Circuit Design}
A variable DC voltage source is connected across the main terminals. A gate current is provided to trigger the device. The voltage is swept from negative to positive values.


\begin{figure}[H]
    \centering
    \includegraphics[width=0.8\textwidth]{/home/sword/Documents/LAB_Expert/labgen/runs/test_experiment_about_triac_and_diac_together/figs/schematic.png}
    \caption{Experimental Circuit Diagram}
\end{figure}


\section{Required Apparatus}
\begin{itemize}

    \item DC Power Supply (Variable) (Quantity: 1)

    \item Device Under Test (Quantity: 1)

    \item Resistor ($1 k\Omega$) (Quantity: 1)

    \item Multimeter (Quantity: 2)

\end{itemize}

\section{Procedure}
\begin{enumerate}

    \item Connect the circuit as per the experimental circuit diagram.

    \item Apply a constant gate current $I_G$.

    \item Vary the supply voltage from -15V to 15V.

    \item Record the voltage and current.

    \item Plot the characteristics.

\end{enumerate}

\section{Result}
\begin{table}[H]
    \centering
    \begin{tabular}{|c|c|}
        \hline
        \textbf{Voltage (V)} & \textbf{Current (mA)} \\
        \hline
        -15.0 & -759.3 \\
        -11.7 & -753.0 \\
        -8.4 & -7.6 \\
        -5.0 & -3484.0 \\
        -1.7 & -1.5 \\
        1.6 & 712.9 \\
        5.0 & 4.3 \\
        8.3 & 748.2 \\
        11.7 & 10.9 \\
        15.0 & 762.8 \\
        \hline
    \end{tabular}
    \caption{Simulated Data (Sample Points)}
\end{table}



\begin{figure}[H]
    \centering
    \includegraphics[width=0.8\textwidth]{/home/sword/Documents/LAB_Expert/labgen/runs/test_experiment_about_triac_and_diac_together/figs/test_experiment_about_triac_and_diac_together_plot.png}
    \caption{ Simulated Test Experiment about triac and diac together Curve }
\end{figure}



\section{Discussion}
The circuit was constructed by connecting the DIAC in series with a variable resistor to the gate of the TRIAC, while a resistor-inductive load was placed in series with the TRIAC main terminals and the AC source. During the observation of the waveforms, the DIAC did not conduct until the voltage across it reached its breakover threshold, at which point it latched on and injected a sharp current pulse into the TRIAC gate. This pulse triggered the TRIAC, causing it to switch from the blocking state to the conducting state. The simulated results showed that the TRIAC conducted in both the positive and negative half-cycles of the AC supply, confirming its bidirectional functionality. As the resistance in the trigger circuit was increased, the breakover time of the DIAC was delayed, resulting in a larger firing angle and a corresponding reduction in the average voltage delivered to the load. The symmetry of the conduction angle in both half-cycles was verified, indicating that the DIAC provided consistent triggering regardless of the AC polarity. Minor deviations in the trigger pulses were observed, which were attributed to component tolerances and the inherent variation in the DIAC breakover voltage.

\section{Conclusion}
The experiment successfully demonstrated the functional interplay between a TRIAC and a DIAC in an AC power control application. It was verified that the DIAC acts as a bidirectional trigger source, initiating conduction in the TRIAC only when its breakover voltage is exceeded. The TRIAC exhibited its characteristic latching behavior, conducting current in both directions until the main terminal current fell below the holding current. By varying the firing angle through the RC timing circuit, the average power transferred to the load was effectively regulated, showcasing the viability of this configuration for dimming and speed control applications. The simulation results confirmed that the DIAC ensures symmetric triggering in both AC half-cycles, preventing the power imbalance that would occur with a standard SCR. Through this test, the underlying principles of phase-controlled AC power regulation were validated, highlighting the importance of semiconductor thyristors in modern power electronics circuits. The observed waveforms and voltage levels aligned with theoretical expectations, confirming the correct operation of the components and the circuit design.

\section{References}
\begin{enumerate}

    \item Generated by LabGen Knowledge Hub API

    \item Ngspice Simulation Data.

\end{enumerate}

\end{document}